\documentclass[11pt,letterpaper]{article}

\usepackage[top=1in, bottom=1in, left=1in, right=1in, headheight=14pt]{geometry}
\usepackage{fontspec}
\setmainfont{DejaVu Serif}
\setsansfont{DejaVu Sans}
\setmonofont{DejaVu Sans Mono}

\usepackage{xcolor}
\usepackage{titlesec}
\usepackage{fancyhdr}
\usepackage{enumitem}
\usepackage{hyperref}
\usepackage{booktabs}
\usepackage{tabularx}
\usepackage{longtable}
\usepackage{colortbl}
\usepackage{multirow}
\usepackage{array}
\usepackage{parskip}
\usepackage{tcolorbox}
\usepackage{graphicx}
\usepackage{lastpage}

% --- Ocean/Maritime Color Scheme ---
\definecolor{primary}{HTML}{0D47A1}
\definecolor{secondary}{HTML}{00695C}
\definecolor{tertiary}{HTML}{5D4037}
\definecolor{accent}{HTML}{1565C0}
\definecolor{lightprimary}{HTML}{E3F2FD}
\definecolor{lightsecondary}{HTML}{E0F2F1}
\definecolor{lighttertiary}{HTML}{EFEBE9}
\definecolor{rowgray}{HTML}{F5F5F5}
\definecolor{bordergray}{HTML}{BDBDBD}
\definecolor{subtlegray}{HTML}{757575}
\definecolor{darktext}{HTML}{212121}

\newcolumntype{L}[1]{>{\raggedright\arraybackslash}p{#1}}
\newcolumntype{C}[1]{>{\centering\arraybackslash}p{#1}}
\newcommand{\tableheader}[1]{\textbf{\sffamily\textcolor{white}{#1}}}

\titleformat{\section}
  {\Large\bfseries\sffamily\color{primary}}
  {\thesection}{1em}{}
  [\vspace{-0.5em}\textcolor{bordergray}{\rule{\textwidth}{0.4pt}}]

\titleformat{\subsection}
  {\large\bfseries\sffamily\color{secondary}}
  {\thesubsection}{1em}{}

\titleformat{\subsubsection}
  {\normalsize\bfseries\sffamily\color{tertiary}}
  {\thesubsubsection}{1em}{}

\titlespacing*{\section}{0pt}{1.5em}{0.8em}
\titlespacing*{\subsection}{0pt}{1.2em}{0.5em}
\titlespacing*{\subsubsection}{0pt}{0.8em}{0.3em}

\newcommand{\stageheader}[2]{%
  \noindent\colorbox{#2}{\makebox[\textwidth][l]{%
    \textcolor{white}{\bfseries\sffamily\hspace{6pt}#1\hspace{6pt}}}}%
  \vspace{0.5em}\par%
}

\newtcolorbox{asciibox}{
  colback=rowgray, colframe=bordergray,
  arc=1pt, boxrule=0.5pt,
  left=6pt, right=6pt, top=4pt, bottom=4pt,
  fontupper=\small\ttfamily,
  breakable
}

\newtcolorbox{quotebox}{
  colback=lightprimary, colframe=primary,
  arc=2pt, boxrule=0.5pt,
  left=12pt, right=12pt, top=8pt, bottom=8pt,
  breakable
}

\newcommand{\metafield}[2]{\textbf{\sffamily #1} #2\par}

\pagestyle{fancy}
\fancyhf{}
\fancyhead[L]{\small\sffamily\textcolor{subtlegray}{Maritime Platform Cost Model}}
\fancyhead[R]{\small\sffamily\textcolor{subtlegray}{GSD Mission Package}}
\fancyfoot[C]{\small\sffamily\textcolor{subtlegray}{Page \thepage\ of \pageref{LastPage}}}
\renewcommand{\headrulewidth}{0.4pt}
\renewcommand{\footrulewidth}{0pt}

\hypersetup{
  colorlinks=true,
  linkcolor=secondary,
  urlcolor=secondary,
  pdfauthor={GSD Ecosystem},
  pdftitle={Maritime Compute + Energy + Transport Platform -- Integrated Cost Model},
  pdfsubject={Vision to Mission Pipeline},
}

\begin{document}

% --- TITLE PAGE ---
\thispagestyle{empty}
\begin{center}
\vspace*{3cm}

{\small\sffamily\textcolor{primary}{\textbf{GSD RESEARCH MISSION PACKAGE}}}

\vspace{0.3cm}
\textcolor{bordergray}{\rule{8cm}{0.4pt}}

\vspace{1.5cm}
{\fontsize{28}{34}\selectfont\bfseries\sffamily\textcolor{primary}{Maritime Platform\\[0.3em]Integrated Cost Model}}

\vspace{0.8cm}
{\Large\sffamily\textcolor{secondary}{Compute + Energy + Transport --- Full Platform Economics}}

\vspace{0.5cm}
{\large\sffamily\itshape\textcolor{tertiary}{A Deep Research Study}}

\vspace{3cm}
{\sffamily\textcolor{accent}{Vision $\rightarrow$ Research $\rightarrow$ Mission}}\\[0.3em]
{\small\sffamily\textcolor{accent}{Full Pipeline Execution}}

\vspace{3cm}
{\small\sffamily\textcolor{subtlegray}{Date: March 26, 2026  |  Status: Ready for GSD Orchestrator}}\\[0.3em]
{\small\sffamily\textcolor{subtlegray}{Activation Profile: Fleet (16 roles)  |  Estimated: 6--8 context windows across 3 sessions}}
\end{center}
\newpage

% --- TABLE OF CONTENTS ---
\tableofcontents
\newpage

% =====================================================================
\stageheader{STAGE 1 --- VISION DOCUMENT}{primary}

\section{Maritime Platform Cost Model --- Vision Guide}

\metafield{Date:}{March 26, 2026}
\metafield{Status:}{Research-Ready / Pre-Execution}
\metafield{Depends On:}{Fox Infrastructure Group Master Business Plan, GSD Mesh Prototype Spec, Deep Audio Analyzer Mission}
\metafield{Context:}{Fox Infrastructure Group, GSD Skill-Creator v1.50 Research Track}

\subsection{Vision}

The spaces between nodes are where the money lives. Individual cost data for a maritime compute platform
exists in scattered form: Norway's Submerged Floating Tunnel concept runs \$25 billion for 3.7
kilometers, subsea power cables span \$0.5--5 million per kilometer depending on voltage and depth,
shore-based data centers average \$10--12 million per megawatt, and a 6,000-TEU container vessel
runs \$80--100 million off the Korean yards. What does not yet exist is a model that shows what
these numbers become when they converge --- when the same cable serves both power delivery and
data transport, when the same vessel hull doubles as a compute chassis chassis and a coolant heat
exchanger, when the port infrastructure serves both cargo and maintenance operations.

This mission answers the gap stated directly: cost analysis across the full platform is absent.
The research module has hit its ceiling on individual component data. The integrated cost of a
maritime compute plus energy plus transport platform requires system-level modeling that treats
convergence credits as first-class cost items. The Amiga Principle, applied to infrastructure
economics, asks the same question it always asks: what does architectural leverage save you?
A platform where compute cooling is provided for free by the surrounding ocean, where power
delivery cables double as data backhaul, where vessels en route between nodes also serve as
mobile compute edge --- that platform's cost is not the sum of its parts. The gap this mission
closes is the model that proves it.

The Fox Infrastructure Group vision spans BC to Chile, with FoxCompute anchored at Electric City
(Grand Coulee / BPA hydro at \$0.02/kWh) and FoxFiber running submarine cable down the coast.
The maritime question is: what does it cost to extend that compute and energy spine offshore and
make it mobile? The answer requires a six-module integrated model --- compute, energy, transport,
convergence credits, scenario comparison, and sensitivity --- executed by a Fleet crew with enough
depth to produce actionable numbers, not ranges.

\subsection{Problem Statement}

\begin{enumerate}
\item \textbf{No integrated platform cost model exists.} Individual component costs are documented
  separately for data centers (\$10--12M/MW), subsea cable (\$0.5--5M/km), offshore wind (\$2,852/kW
  installed), and container vessels (\$15--200M). No published model computes their convergence
  into a single maritime platform CAPEX/OPEX envelope.

\item \textbf{Convergence credits are unmeasured.} A maritime platform shares infrastructure across
  compute, energy, and transport functions (ocean cooling, cable dual-use, vessel co-location).
  These shared-infrastructure savings --- potentially 20--40\% on individual component costs ---
  have not been quantified for the platform as a whole.

\item \textbf{Scenario range is undefined.} The platform could be deployed at prototype scale (1--5 MW),
  regional scale (50--150 MW), or full-corridor scale (500 MW+). The cost structure changes
  dramatically across these scales, but no comparison framework exists.

\item \textbf{Operating cost dominates lifetime economics.} Subsea maintenance, vessel day rates,
  cable repair costs, and offshore power opex can dwarf initial CAPEX over a 20-year platform
  life. Current planning has focused on CAPEX without an opex model.

\item \textbf{Risk-adjusted returns are unmodeled.} Offshore infrastructure faces weather, geopolitical,
  and regulatory risks not present in land-based systems. No risk-adjusted NPV analysis has
  been produced for the platform at any deployment scale.
\end{enumerate}

\subsection{Core Concept}

\textbf{One-line model:} A maritime compute-energy-transport platform is an infrastructure stack
where shared physical media (cables, hulls, ocean) reduces integrated system cost below the
arithmetic sum of standalone components, and the correct unit of analysis is the platform
CAPEX/OPEX envelope, not the component cost sheet.

\subsubsection{Platform Architecture Map}

\begin{asciibox}
MARITIME PLATFORM COST ARCHITECTURE

     OFFSHORE GENERATION         SUBSEA SPINE          SHORE ANCHOR
     +-----------------+     +---------------+     +----------------+
     | Floating Wind   |-----| Power + Data  |-----| Electric City  |
     | $2,852/kW (fix) |     | Cable Bundle  |     | $0.02/kWh BPA  |
     | $4,000+/kW (flt)|     | $0.5-5M/km    |     | Grand Coulee   |
     |                 |     |               |     |                |
     | Tidal/Wave      |     | Dual-use:     |     | Columbia River |
     | (emerging)      |     |  - Power HVDC |     | Free cooling   |
     +-----------------+     |  - Data fiber |     +----------------+
                             +-------+-------+
                                     |
              +-----------------+    |    +-----------------+
              | Compute Node A  |----+----| Compute Node B  |
              | Container DC    |         | Vessel-mounted  |
              | $10-12M/MW land |         | (marine premium |
              | +30-50% marine  |         |  offset by      |
              |                 |         |  free cooling)  |
              +-----------------+         +-----------------+
                                     |
                             +-------+-------+
                             | Mobile Edge   |
                             | Cargo vessels |
                             | en-route      |
                             | compute nodes |
                             +---------------+

CONVERGENCE CREDIT OPPORTUNITIES:
  C1: Ocean cooling eliminates chiller CAPEX (~$1-2M/MW saved)
  C2: Cable bundle dual-use (power + data) vs. separate runs
  C3: Vessel hull compute: CapEx shared with cargo function
  C4: Port infrastructure: maintenance + logistics shared
  C5: BPA hydro anchor: $0.02/kWh vs $0.12 US average
\end{asciibox}

\subsubsection{Research Module Architecture}

\begin{asciibox}
WAVE STRUCTURE --- 6 RESEARCH MODULES

Wave 0 (Foundation):
  F-0: Shared schemas, cost taxonomy, source index

Wave 1A (Parallel Track):              Wave 1B (Parallel Track):
  MOD-A: Compute Cost Analysis           MOD-B: Energy Infrastructure Costs
  - Offshore DC economics                - Floating wind: CAPEX + LCOE
  - Marine premium vs. cooling credit    - Subsea cable: cost per km
  - Container DC benchmarks              - BPA hydro anchor economics
  - AI facility requirements             - Grid interconnect costs

Wave 1C (Parallel Track):              Wave 1D (Parallel Track):
  MOD-C: Maritime Transport Economics    MOD-D: Shared Infrastructure Credits
  - Vessel acquisition/operating         - Convergence credit model
  - Autonomous shipping economics        - Dual-use infrastructure savings
  - Port infrastructure costs            - Ocean cooling quantification
  - Maintenance vessel day rates         - Shared opex reduction

Wave 2 (Synthesis):
  MOD-E: Scenario Comparison (3 scales: 5MW / 100MW / 500MW)
  MOD-F: Sensitivity & Risk Analysis (NPV, risk-adjusted returns)

Wave 3 (Publication + Verification):
  Integration audit, source validation, PDF assembly
\end{asciibox}

\subsection{Research Modules}

\subsubsection{MOD-A: Compute Cost Analysis}
Offshore and marine-deployed data center economics. Baseline shore-based construction cost
(\$10--12M/MW general, \$20M+/MW AI-optimized), marine deployment premium (estimated 30--50\%
structural and weatherproofing uplift), offset by ocean cooling savings (elimination of
chiller plant, estimated \$1--2M/MW). Container DC construction benchmarks. AI vs.\ standard
compute density implications for the platform.

\subsubsection{MOD-B: Energy Infrastructure Costs}
Floating offshore wind CAPEX (\$4,000+/kW 2024, targeting \$45/MWh LCOE by 2035 per DOE),
fixed offshore wind (\$2,852/kW IRENA 2024), subsea power cable costs (\$0.5--5M/km depending
on voltage, depth, and seabed conditions), HVAC vs.\ HVDC break-even distance (\textasciitilde{}80km),
BPA hydro baseline (\$0.02/kWh Grand Coulee), grid interconnect infrastructure costs.

\subsubsection{MOD-C: Maritime Transport Economics}
Vessel acquisition cost by scale (\$15--50M small, \$80--200M large), annual operating costs
(crew \$2--5M/yr, fuel \$3--10M/yr), autonomous vessel economics (3.4\% freight rate reduction
per ScienceDirect 2017, shore control center amortization), port infrastructure investment,
maintenance vessel day rates (\$15,000--50,000/day offshore), lifecycle and drydock costs.

\subsubsection{MOD-D: Shared Infrastructure Credit Model}
Quantification of convergence credits: (C1) ocean cooling CAPEX elimination, (C2) cable
bundle dual-use savings vs.\ separate power and fiber runs, (C3) vessel hull compute
co-location economics, (C4) shared port maintenance and logistics infrastructure,
(C5) BPA hydro power cost advantage over grid average. Net convergence credit as
percentage reduction on integrated platform cost vs.\ standalone components.

\subsubsection{MOD-E: Scenario Comparison}
Three deployment scales with full CAPEX/OPEX/NPV analysis: (1) Prototype --- 5 MW compute,
1 vessel, 50km cable run; (2) Regional --- 100 MW compute, 5 vessels, 300km spine;
(3) Full Corridor --- 500 MW compute, 20+ vessels, BC-to-Chile spine integration
with FoxFiber. Cost per MW-hour of compute delivered as primary comparison metric.

\subsubsection{MOD-F: Sensitivity and Risk Analysis}
Key cost drivers and their sensitivity ranges: steel prices (20--30\% swing on vessel/structure
costs), cable installation vessel day rates (\$150,000--300,000/day), offshore wind capex
trajectory (40--60\% projected decline by 2035 per IRENA), regulatory and permitting risk
premiums, weather event downtime modeling. Risk-adjusted NPV across three scenarios.
Insurance and bond cost requirements for maritime infrastructure.

\subsection{Chipset Configuration}

\begin{asciibox}
chipset:
  mission_id: "maritime-cost-model-v1.0"
  profile: "fleet"
  modules:
    - id: "MOD-A"
      name: "Compute Cost Analysis"
      model: "opus"
      role: "CRAFT-COMPUTE"
      parallel_group: "wave_1a"
    - id: "MOD-B"
      name: "Energy Infrastructure Costs"
      model: "sonnet"
      role: "CRAFT-ENERGY"
      parallel_group: "wave_1a"
    - id: "MOD-C"
      name: "Maritime Transport Economics"
      model: "sonnet"
      role: "CRAFT-MARITIME"
      parallel_group: "wave_1b"
    - id: "MOD-D"
      name: "Shared Infrastructure Credits"
      model: "opus"
      role: "CRAFT-CONVERGENCE"
      parallel_group: "wave_1b"
    - id: "MOD-E"
      name: "Scenario Comparison"
      model: "opus"
      role: "ANALYST"
      parallel_group: "wave_2"
      depends_on: ["MOD-A", "MOD-B", "MOD-C", "MOD-D"]
    - id: "MOD-F"
      name: "Sensitivity and Risk"
      model: "opus"
      role: "ANALYST"
      parallel_group: "wave_2"
      depends_on: ["MOD-A", "MOD-B", "MOD-C", "MOD-D"]
  foundation:
    - id: "F-0"
      name: "Cost Taxonomy and Source Index"
      model: "haiku"
      role: "LOG"
      parallel_group: "wave_0"
\end{asciibox}

\subsection{Scope Boundaries}

\subsubsection{In Scope (v1.0)}
\begin{itemize}[itemsep=2pt]
  \item CAPEX and OPEX modeling for offshore compute, subsea power/data cable, maritime vessels
  \item Three deployment scales: 5 MW prototype, 100 MW regional, 500 MW corridor
  \item Convergence credit quantification for shared infrastructure
  \item 20-year NPV analysis with risk adjustment
  \item BPA hydro power anchor as cost baseline reference
  \item Autonomous vs.\ manned vessel economics comparison
  \item Floating vs.\ fixed offshore wind cost comparison at the platform scale
\end{itemize}

\subsubsection{Out of Scope}
\begin{itemize}[itemsep=2pt]
  \item Specific site permitting costs (jurisdiction-dependent, treated as risk range)
  \item Detailed naval architecture or hull engineering specifications
  \item Customs, duties, and cross-border regulatory costs by individual country
  \item Satellite-based data transport costs (assumed fiber cable primary)
  \item Carbon credit or environmental offset revenue modeling
  \item Full financial model (investment structure, equity/debt splits, tax treatment)
\end{itemize}

\subsection{Success Criteria}

\begin{enumerate}
  \item Integrated platform CAPEX produced for all three scales (5 MW, 100 MW, 500 MW)
    with component-level breakdown and source citation for every line item.
  \item Integrated platform OPEX (annual) produced for all three scales with
    maintenance, fuel/power, crew, and cable repair costs itemized.
  \item Convergence credit total quantified for each scale as a percentage reduction
    from standalone component sum, with methodology documented.
  \item Cost per MW-hour of compute delivered calculated for each scenario as the
    primary platform economics comparison metric.
  \item Offshore wind CAPEX trajectory modeled from 2024 to 2035, with DOE \$45/MWh
    LCOE target year and implied platform cost reduction documented.
  \item Subsea cable cost per km compiled across voltage class, depth band, and
    seabed condition with at least five sourced real-project benchmarks.
  \item Vessel economics compared between manned and autonomous operation at
    platform scale with sourced operating cost assumptions.
  \item Ocean cooling credit quantified in \$/MW saved vs.\ shore-based chiller plant,
    with engineering basis documented.
  \item 20-year NPV analysis produced for the 100 MW regional scenario at three
    discount rates (5\%, 8\%, 12\%).
  \item Sensitivity tornado chart produced identifying the top five cost drivers
    and their percentage impact on integrated platform cost.
  \item All numerical claims attributed to specific sources (government agencies,
    peer-reviewed journals, industry reports from named firms).
  \item Research Reference document is self-contained --- a reader with no prior
    context can reconstruct the cost model from the document alone.
\end{enumerate}

\subsection{Relationship to Other Documents}

\begin{center}
\rowcolors{2}{rowgray}{white}
\begin{tabularx}{\textwidth}{L{4.5cm}L{3cm}X}
\toprule
\rowcolor{primary}
\tableheader{Document} & \tableheader{Relationship} & \tableheader{Dependency} \\
\midrule
Fox Infrastructure Group Master Plan & Parent & Electric City BPA anchor, 21-company structure, corridor geography \\
GSD Mesh Prototype Spec & Sibling & 10-node white-box cluster; land-based compute cost baseline \\
Fox Legal Staffing Analysis & Sibling & Regulatory and permitting cost inputs for maritime operations \\
Deep Audio Analyzer Mission & Sibling & Example platform mission structure; execution model reference \\
GSD Skill-Creator v1.50 & Consumer & Cost model outputs feed infrastructure planning wave \\
\bottomrule
\end{tabularx}
\end{center}

\subsection{The Through-Line}

The Amiga Principle does not stop at software architecture. The same insight that built
the Amiga --- that specialized, purpose-built components working in concert produce outcomes
impossible for generalized brute-force systems --- applies to physical infrastructure at
maritime scale. A shore-based data center optimized for compute but paying \$0.12/kWh and
running chiller plants is a generalized system. A maritime platform where the ocean provides
cooling, the cable provides both power and backhaul, the vessels provide both compute and
logistics, and the BPA dam provides baseload at \$0.02/kWh is an Amiga. The question this
mission answers is not whether the architecture is elegant. That was settled in the vision.
The question is what it costs to build the Amiga at sea --- and whether the convergence
credits are large enough to close the gap between the marine deployment premium and the
infrastructure savings. That is not a philosophical question. That is a number. This mission
produces it.

% =====================================================================
\newpage
\stageheader{STAGE 2 --- RESEARCH REFERENCE}{secondary}

\section{Maritime Platform Cost Model --- Research Reference}

\subsection{How to Use This Document}

This reference compiles sourced cost data organized by module. All monetary values are
2024 USD unless otherwise noted. All claims are attributed to specific sources. Use
this document as the primary data input for the mission execution wave plan.

\textbf{Key source organizations:}
\begin{itemize}[itemsep=2pt]
  \item International Renewable Energy Agency (IRENA) --- offshore wind cost database
  \item U.S. Department of Energy / NREL --- Annual Technology Baseline (ATB), Cost of Wind Energy Review
  \item Wood Mackenzie --- global LCOE reports (2024)
  \item Cushman \& Wakefield --- Data Center Development Cost Guide (2025)
  \item Norwegian Public Roads Administration (NPRA) --- subsea tunnel and cable infrastructure
  \item SINTEF / NorthWind --- floating wind cost analysis (Norway, 2024)
  \item OilPrice.com / IEEE Spectrum --- subsea cable cost benchmarks
  \item IMO / MDPI Applied Sciences --- autonomous vessel economics
  \item CBRE / Dgtl Infra --- data center cost benchmarks
\end{itemize}

\subsection{MOD-A Reference: Compute Cost Analysis}

\subsubsection{Shore-Based Data Center Benchmarks (2024--2025)}

\begin{center}
\rowcolors{2}{rowgray}{white}
\begin{tabularx}{\textwidth}{L{4.5cm}C{2.5cm}C{2.5cm}L{3cm}}
\toprule
\rowcolor{primary}
\tableheader{Facility Type} & \tableheader{CAPEX/MW} & \tableheader{CAPEX/sqft} & \tableheader{Source} \\
\midrule
General purpose (Tier III) & \$10--12M & \$600--1,100 & CBRE / Cushman 2025 \\
AI-optimized (liquid cool) & \$20M+ & \$1,200--2,200 & CBRE 2025 \\
Hyperscale campus & \$12M avg & \$1,050 & Digital Realty 2024 \\
Brownfield redevelopment & \$7--8M & \$525--700 & QTS / Dgtl Infra 2024 \\
Prime market (NYC/SV) & \$18--25M & \$2,200 & Cushman 2024 \\
\bottomrule
\end{tabularx}
\end{center}

\textbf{Key cost components:} MEP systems consume 40--50\% of total data center budget. Electrical
systems alone represent 40--45\% of expenditures. Land represents 15--20\% of total cost at
\$5.59/sqft average in 2024, with large-parcel land costs up 23\% year-over-year (Cushman
\& Wakefield 2025).

\subsubsection{Marine Deployment Premium and Ocean Cooling Credit}

Shore-based data centers do not apply to maritime deployment without three adjustments:

\textbf{Marine Structural Premium (estimated +30--50\% on shell):}
\begin{itemize}[itemsep=2pt]
  \item Corrosion-resistant materials and coatings for saltwater environments
  \item Structural reinforcement for wave loading and vessel motion (if vessel-mounted)
  \item Weatherproofing and pressure-rated enclosures for subsurface installations
  \item Marine electrical standards (IEC 60092) vs.\ land-based codes
\end{itemize}

\textbf{Ocean Cooling Credit (estimated --\$1--2M/MW CAPEX, --\$0.5--1M/MW annual OPEX):}
\begin{itemize}[itemsep=2pt]
  \item Shore-based chiller plant accounts for \$1--2M/MW of MEP CAPEX (derived: MEP = 50\% of
    \$10M/MW = \$5M/MW; cooling = 20--40\% of MEP = \$1--2M/MW)
  \item North Atlantic and North Pacific surface water temperatures 8--15°C provide
    direct cooling without mechanical refrigeration
  \item Seawater cooling eliminates 15--20\% of ongoing electricity consumption
    attributable to cooling systems in shore-based facilities
  \item Water availability concern (shore-based facilities consume \textasciitilde{}2M liters/day per
    100 MW per IEA 2024) is eliminated: ocean provides free, unlimited supply
\end{itemize}

\textbf{Net marine compute CAPEX estimate:} \$10M/MW base $\times$ 1.35 marine premium $-$ \$1.5M/MW
cooling credit $=$ \$12M/MW. Comparable to AI-optimized shore-based at equivalent compute
density, with superior cooling headroom for future density growth.

\subsection{MOD-B Reference: Energy Infrastructure Costs}

\subsubsection{Offshore Wind: CAPEX and LCOE (2024)}

\begin{center}
\rowcolors{2}{rowgray}{white}
\begin{tabularx}{\textwidth}{L{4cm}C{2.5cm}C{2.5cm}X}
\toprule
\rowcolor{secondary}
\tableheader{Technology} & \tableheader{CAPEX/kW} & \tableheader{LCOE} & \tableheader{Source} \\
\midrule
Fixed-bottom offshore wind & \$2,852/kW & \$79/MWh & IRENA 2024 \\
Fixed offshore (Wood Mac) & --- & \$230/MWh & Wood Mackenzie 2024 \\
Floating offshore (current) & \$4,000+/kW & \$320/MWh & Wood Mackenzie 2024 \\
Floating (WindFloat pilot) & --- & \$112/MWh & MDPI FOW Review 2024 \\
DOE Floating Wind Shot target & --- & \$45/MWh & DOE NREL 2024 (by 2035) \\
ETIPWind target at 20 GW & --- & €40--80/MWh & SINTEF/NorthWind 2024 \\
Onshore wind (benchmark) & --- & \$75/MWh & Wood Mackenzie 2024 \\
BPA Grand Coulee hydro & --- & \$0.02/kWh & Fox Infrastructure Group \\
\bottomrule
\end{tabularx}
\end{center}

\textbf{Capacity factors:} Modern fixed-bottom offshore reaches 40--60\% capacity factor
(IRENA 2024). Floating offshore currently lower due to immaturity; projected to converge
with fixed-bottom by 2030--2035 as supply chains mature.

\textbf{CAPEX trajectory:} Global weighted average installed cost for offshore wind has fallen
62\% between 2010 and 2024 (IRENA 2024). Total installed offshore wind capacity grew from
3.1 GW in 2010 to 82.9 GW in 2024. Chinese suppliers offering turbines at \$300,000/MW (SINTEF
Oct.\ 2024) vs.\ European land-based at \$1,400/MW --- demonstrating the scale of supply chain
maturity impact.

\subsubsection{Subsea Power Cable Costs}

\begin{center}
\rowcolors{2}{rowgray}{white}
\begin{tabularx}{\textwidth}{L{4cm}C{3cm}X}
\toprule
\rowcolor{secondary}
\tableheader{Cable Type} & \tableheader{Installed Cost/km} & \tableheader{Notes / Source} \\
\midrule
HVAC medium voltage (\textless{}150 kV) & \$0.5--2.0M/km & Short, shallow, protected routes (Quora industry 2024) \\
HVDC (\textpm{}100--300 kV) & \$1.0--4.0M/km & Standard interconnectors (Quora 2024) \\
Complex / erratic seabed & \$3.0--5.0M/km & Deep water, rocky terrain (Thunder Said Energy) \\
Offshore wind inter-array & \$0.5M/km & Within-farm cabling (Thunder Said Energy) \\
Subsea power cable (general) & \$2.5M+/km & OilPrice.com 2023 \\
EuroAsia Interconnector & \$0.74M/km & 1,208 km at \$900M (Submarine Networks 2021) \\
\bottomrule
\end{tabularx}
\end{center}

\textbf{AC vs.\ HVDC:} AC cables are economical for distances under \textasciitilde{}80--100 km; beyond
that, HVDC becomes cost-effective despite higher converter station costs (ESRU/Strathclyde).
For a BC-to-Chile corridor spanning 10,000+ km, HVDC is the only viable technology.

\textbf{Installation vessels:} Cable installation vessels cost \$150,000--300,000/day to charter.
State-of-the-art vessels such as Prysmian's Leonardo da Vinci cost \$200M to build and can
lay 2.1 km/hour in depths up to 3,000 m (Thunder Said Energy). Installation is often the
cost-controlling factor in complex routes.

\textbf{Communication fiber context:} Submarine communication cables cost \$30,000--50,000/km
(OilPrice.com 2023 / Thematica) --- two orders of magnitude cheaper than power cables.
Bundling power and fiber in the same cable trench provides installation cost savings even
if the cables themselves are separate systems.

\subsection{MOD-C Reference: Maritime Transport Economics}

\subsubsection{Vessel Acquisition Costs (2024--2026)}

\begin{center}
\rowcolors{2}{rowgray}{white}
\begin{tabularx}{\textwidth}{L{4cm}C{3cm}X}
\toprule
\rowcolor{tertiary}
\tableheader{Vessel Type / Size} & \tableheader{New Build Cost} & \tableheader{Source} \\
\midrule
Small cargo (1,000--5,000 TEU) & \$15--50M & FreightAmigo / Bearcat 2026 \\
Standard freight (5,000--10,000 TEU) & \$30--100M & FreightAmigo 2024 \\
Large container (10,000+ TEU) & \$100--200M & FreightAmigo / Wikipedia 2024 \\
Ultra-large (20,000+ TEU) & \$200--267M & Accio.com market data 2024 \\
4,000--6,000 TEU (Korean yards) & \$80--100M & Industry operator quoted 2024 \\
Used vessel (5-year-old) & 60--70\% of new & Bearcat Express 2026 \\
\bottomrule
\end{tabularx}
\end{center}

\subsubsection{Annual Operating Costs (Large Vessel, indicative)}

\begin{center}
\rowcolors{2}{rowgray}{white}
\begin{tabularx}{\textwidth}{L{4cm}C{3cm}X}
\toprule
\rowcolor{tertiary}
\tableheader{Cost Item} & \tableheader{Annual Range} & \tableheader{Notes} \\
\midrule
Crew wages + related & \$2--5M & Route and flag-state dependent \\
Fuel (large vessel) & \$3--10M & Route-dependent; dominant variable cost \\
IMO 2020 sulfur surcharge & \$150--300/container & \textasciitilde{}\$500K--1.5M/yr at 3,000 TEU \\
EU ETS (European ports) & \$50--150/container & From 2024 phase-in \\
Annual maintenance/drydock & \$1--3M (amortized) & Drydock every 3--5 years \\
Insurance & 0.5--1.5\% of vessel value & \\
Port fees (per voyage) & Variable & \\
\bottomrule
\end{tabularx}
\end{center}

\subsubsection{Autonomous Vessel Economics}

Research from ScienceDirect (2017, Gdynia Maritime University literature review) and
follow-on MUNIN / ReVolt / Yara Birkeland program data shows:
\begin{itemize}[itemsep=2pt]
  \item Autonomous operation expected present value savings of \$4.3M over a 25-year bulk carrier lifecycle
  \item Required freight rate 3.4\% lower for equivalent autonomous bulker vs.\ conventional manned ship
  \item Shore Control Centre can monitor approximately 90 vessels simultaneously (amortizes SCC cost)
  \item Communication costs increase \textasciitilde{}10$\times$ vs.\ conventional vessel (significant for data-heavy compute payloads)
  \item Yara Birkeland: 120 TEU, fully electric, fully autonomous coastal feeder --- production unit economics pending 2025--2026 operational data
\end{itemize}

\textbf{Platform implication:} For compute-carrying vessels that are also running
workloads en route, the 10$\times$ communication cost premium is offset by the compute
revenue those workloads generate. This is the central economic argument for vessel-mounted
edge compute as a maritime platform component.

\subsection{MOD-D Reference: Shared Infrastructure Credit Model}

\textbf{Convergence Credits --- Quantification Basis}

\begin{center}
\rowcolors{2}{rowgray}{white}
\begin{tabularx}{\textwidth}{L{1cm}L{5cm}C{2.5cm}X}
\toprule
\rowcolor{primary}
\tableheader{ID} & \tableheader{Credit Type} & \tableheader{Estimated Value} & \tableheader{Basis} \\
\midrule
C1 & Ocean cooling eliminates chiller CAPEX & \$1--2M/MW & Cooling = 20--40\% of MEP = 20--40\% of 50\% of \$10M/MW \\
C2 & Cable bundle installation: power + fiber in same trench & 20--30\% on cable install & Single mobilization of \$150K--300K/day vessel serves both \\
C3 & Vessel hull compute co-location vs.\ standalone platform & 40--60\% on float structure & Hull already amortized by cargo function \\
C4 & Shared port logistics (cargo + maintenance) & 15--25\% on port OPEX & Berth, crane, and service infrastructure shared \\
C5 & BPA hydro power vs.\ grid average (\$0.02 vs.\ \$0.12/kWh) & \$438M/yr at 500 MW & Fox Infrastructure Group electric city model \\
\bottomrule
\end{tabularx}
\end{center}

\textbf{Net convergence credit estimate:} At 100 MW regional scale, credits C1--C4 are
estimated to reduce integrated CAPEX by 18--28\% vs.\ standalone component sum. Credit C5
(BPA power) applies only to shore-anchored compute and reduces annual OPEX by \$87M vs.\
US average at 500 MW scale. The largest single credit (C5) is therefore a land-side
advantage that makes the maritime platform's power economics depend heavily on the cable
cost and loss profile connecting to the BPA anchor.

\subsection{MOD-E Reference: Scenario Comparison Data}

\subsubsection{Prototype Scale (5 MW Compute, 1 Vessel, 50 km Cable)}

\begin{center}
\rowcolors{2}{rowgray}{white}
\begin{tabularx}{\textwidth}{L{5cm}C{3cm}X}
\toprule
\rowcolor{primary}
\tableheader{Cost Item} & \tableheader{Estimate} & \tableheader{Basis} \\
\midrule
Compute infrastructure & \$60M & 5 MW $\times$ \$12M/MW (marine premium net of cooling) \\
Offshore wind generation & \$14M & 5 MW $\times$ \$2,852/kW (fixed) or \$20M floating \\
Subsea cable (50 km) & \$50--100M & \$1--2M/km HVAC; low end for protected route \\
Platform vessel & \$20--50M & Small modified cargo/industrial vessel \\
Shore interconnect & \$5--10M & Grid tie and converter infrastructure \\
\textbf{CAPEX Total (est.)} & \textbf{\$149--244M} & \\
Annual OPEX & \$8--15M & Cable maintenance, vessel ops, power, staffing \\
\bottomrule
\end{tabularx}
\end{center}

\subsubsection{Regional Scale (100 MW Compute, 5 Vessels, 300 km Spine)}

\begin{center}
\rowcolors{2}{rowgray}{white}
\begin{tabularx}{\textwidth}{L{5cm}C{3cm}X}
\toprule
\rowcolor{secondary}
\tableheader{Cost Item} & \tableheader{Estimate} & \tableheader{Basis} \\
\midrule
Compute infrastructure & \$1.2B & 100 MW $\times$ \$12M/MW \\
Offshore wind (100 MW) & \$285M & 100 MW $\times$ \$2,852/kW \\
Subsea power cable (300 km) & \$300--600M & \$1--2M/km HVDC + converter stations \\
Subsea fiber bundle & \$10--15M & \$30--50K/km co-installed \\
Platform vessels (5) & \$100--500M & \$20M--100M each depending on type \\
Port infrastructure & \$50--100M & Dedicated berth and maintenance facility \\
Shore interconnect + BPA tie & \$30--50M & Grid upgrade and converter \\
\textbf{CAPEX Total (est.)} & \textbf{\$1.97--2.75B} & \\
Annual OPEX & \$80--150M & Power (\$17.5M at \$0.02/kWh), cable, vessels, staff \\
Convergence credit (18--28\%) & --\$355--770M & Applied to CAPEX total \\
\textbf{Net CAPEX after credits} & \textbf{\$1.2--2.4B} & \\
\bottomrule
\end{tabularx}
\end{center}

\subsubsection{Full Corridor Scale (500 MW Compute, 20+ Vessels, BC--Chile Spine)}

At full corridor scale, the platform integrates with the Fox Infrastructure Group
BC-to-Chile spine. Cable runs span 10,000+ km (HVDC required), compute nodes are
distributed at major port intervals, and vessel fleet provides both logistics and
mobile edge compute. Full corridor CAPEX is estimated at \$8--15B, with BPA hydro
power economics (\$0.02/kWh) providing \$438M/yr in power cost savings vs.\ US
grid average at 500 MW scale. This scale requires the detailed NPV analysis
(MOD-F) to assess viability against a 20-year investment horizon.

\subsection{Safety and Sensitivity Considerations}

\begin{center}
\rowcolors{2}{rowgray}{white}
\begin{tabularx}{\textwidth}{L{4cm}L{2.5cm}X}
\toprule
\rowcolor{primary}
\tableheader{Area} & \tableheader{Type} & \tableheader{Consideration} \\
\midrule
Indigenous sovereignty & ABSOLUTE & All cable and platform routes through tribal territorial waters require nation-specific consent and revenue sharing per UNDRIP Article 31 and OCAP® principles. Never generic ``Indigenous'' reference. \\
Source quality & ABSOLUTE & All cost data from named organizations (IRENA, NREL, Wood Mackenzie, CBRE). No blog aggregator data as primary source. \\
Numerical attribution & ABSOLUTE & Every dollar figure attributed to specific named source with date. \\
Marine environment & GATE & Cable and platform installation environmental impact assessments required; no cost data interpreted as environmental clearance. \\
Cost projection uncertainty & ANNOTATE & All forward-looking cost data (especially floating wind 2035 projections) explicitly labeled as projections with named source and scenario. \\
\bottomrule
\end{tabularx}
\end{center}

\subsection{Source Bibliography}

\textbf{Government and Agency Sources}
\begin{itemize}[itemsep=2pt]
  \item International Renewable Energy Agency (IRENA). \textit{Renewable Power Generation Costs in 2024.} IRENA, Abu Dhabi, 2024.
  \item U.S. Department of Energy / NREL. \textit{Cost of Wind Energy Review: 2024 Edition.} NREL/TP-5000-91775. 2025.
  \item NREL. \textit{Annual Technology Baseline (ATB): Offshore Wind 2024.} atb.nrel.gov, 2024.
  \item Norwegian Public Roads Administration (NPRA). \textit{E39 Coastal Highway Project --- Subsea Infrastructure Studies.} 2019--2024.
  \item IEA. \textit{Energy and AI.} International Energy Agency, April 2025.
\end{itemize}

\textbf{Industry and Research Sources}
\begin{itemize}[itemsep=2pt]
  \item Wood Mackenzie. \textit{Global LCOE Reports: Wind and Solar 2024.} WoodMac, October 2024.
  \item Cushman \& Wakefield. \textit{U.S. Data Center Development Cost Guide 2025.} C\&W Research, 2025.
  \item CBRE. \textit{Data Center Construction Cost Benchmarks 2025.} CBRE Research.
  \item Dgtl Infra. ``How Much Does It Cost to Build a Data Center?'' dgtlinfra.com, June 2024.
  \item Thunder Said Energy. \textit{Database of Cable Installation Vessels, Costs.} TSE, 2023.
  \item OilPrice.com. ``Subsea Power Cables: The Future of Global Energy Transport.'' December 2023.
  \item SINTEF / NorthWind. ``Floating Wind Can Be Cheaper Than Expected.'' SINTEF Blog, January 2025.
  \item ENR. ``Norway Sets \$900M Awards for World Record Subsea Road Tunnel.'' January 2023.
\end{itemize}

\textbf{Peer-Reviewed Research}
\begin{itemize}[itemsep=2pt]
  \item Kretschmann et al.\ (2017). ``Analyzing the Economic Benefit of Unmanned Autonomous Ships.'' \textit{Transportation Research Part D.} ScienceDirect.
  \item MDPI. ``Costs and Benefits of Autonomous Shipping --- A Literature Review.'' \textit{Applied Sciences} 11(10), 2021.
  \item MDPI. ``Development and Analysis of a Global Floating Wind Levelised Cost of Energy Map.'' \textit{Clean Technologies} 6(3), September 2024.
\end{itemize}

% =====================================================================
\newpage
\stageheader{STAGE 3 --- MISSION PACKAGE}{tertiary}

\section{Milestone Specification}

\subsection{Mission Objective}

Produce an integrated CAPEX/OPEX cost model for a maritime compute-energy-transport
platform at three deployment scales (5 MW, 100 MW, 500 MW), quantifying convergence
credits from shared infrastructure and producing a 20-year NPV analysis with sensitivity
analysis for the regional (100 MW) scenario. All outputs sourced to professional
organizations or peer-reviewed research; all cost claims numerically attributed.

\subsection{Architecture Overview}

\begin{asciibox}
WAVE ARCHITECTURE --- MARITIME COST MODEL MISSION

  Session 1                Session 2                Session 3
  +------------------+     +------------------+     +------------------+
  | Wave 0           |     | Wave 1 (Parallel)|     | Wave 2 + 3       |
  | F-0: Taxonomy    |---> | MOD-A Compute    |     | MOD-E Scenarios  |
  | Cost schema      |     | MOD-B Energy     |---> | MOD-F Sensitivity|
  | Source index     |     | (Track A)        |     | Integration audit|
  | [Haiku]          |     |                  |     | Publication PDF  |
  +------------------+     | MOD-C Transport  |     | [Opus]           |
                           | MOD-D Convergence|     +------------------+
                           | (Track B)        |
                           | [Sonnet+Opus]    |
                           +------------------+

  Cache TTL: Wave 0 output shared across Wave 1 (5-min window)
  Parallel tracks: MOD-A/B and MOD-C/D run simultaneously
  Gate: CAPCOM approval between Wave 1 and Wave 2
\end{asciibox}

\subsection{System Layers}
\begin{enumerate}
  \item \textbf{Foundation Layer} --- Shared cost taxonomy, unit definitions, source index
  \item \textbf{Component Survey Layer} --- Four parallel modules covering compute, energy, transport, and convergence
  \item \textbf{Integration Layer} --- Scenario comparison and sensitivity analysis
  \item \textbf{Publication Layer} --- Research Reference document, verification, final PDF
\end{enumerate}

\subsection{Deliverables}

\begin{center}
\rowcolors{2}{rowgray}{white}
\begin{tabularx}{\textwidth}{C{0.5cm}L{4.5cm}L{4cm}L{4cm}}
\toprule
\rowcolor{tertiary}
\tableheader{\#} & \tableheader{Deliverable} & \tableheader{Acceptance Criteria} & \tableheader{Spec} \\
\midrule
1 & Cost Taxonomy Schema & All units defined, currency/year baseline set, 6 modules covered & F-0 output, Haiku \\
2 & Compute Cost Analysis (MOD-A) & CAPEX/OPEX table, marine premium, cooling credit quantified & MOD-A document, Opus \\
3 & Energy Infrastructure Report (MOD-B) & Wind CAPEX/LCOE, cable cost/km by class, BPA baseline & MOD-B document, Sonnet \\
4 & Transport Economics Report (MOD-C) & Vessel costs, autonomous economics, port infrastructure & MOD-C document, Sonnet \\
5 & Convergence Credit Model (MOD-D) & 5 credit types quantified, net \% reduction calculated & MOD-D model, Opus \\
6 & Scenario Comparison (MOD-E) & 3-scale CAPEX/OPEX/NPV table, cost/MW-hr metric & MOD-E analysis, Opus \\
7 & Sensitivity Analysis (MOD-F) & Tornado chart top-5 drivers, 20-yr NPV at 3 rates & MOD-F analysis, Opus \\
8 & Research Reference PDF & Complete sourced document, self-contained & Final PDF, Sonnet \\
\bottomrule
\end{tabularx}
\end{center}

\subsection{Component Breakdown}

\begin{center}
\rowcolors{2}{rowgray}{white}
\begin{tabularx}{\textwidth}{L{2.5cm}L{3cm}C{1.5cm}C{2cm}X}
\toprule
\rowcolor{tertiary}
\tableheader{Component} & \tableheader{Dependencies} & \tableheader{Model} & \tableheader{Est. Tokens} & \tableheader{Output} \\
\midrule
F-0 Foundation & None & Haiku & 8K & Cost schema, source index \\
MOD-A Compute & F-0 & Opus & 25K & Marine DC cost analysis \\
MOD-B Energy & F-0 & Sonnet & 30K & Wind, cable, hydro data \\
MOD-C Transport & F-0 & Sonnet & 20K & Vessel, port economics \\
MOD-D Convergence & F-0 & Opus & 20K & Credit quantification \\
MOD-E Scenarios & A,B,C,D & Opus & 35K & 3-scale CAPEX/OPEX/NPV \\
MOD-F Sensitivity & A,B,C,D & Opus & 25K & Tornado, risk-adj NPV \\
Integration Audit & E,F & Sonnet & 15K & Cross-check report \\
Publication & All & Sonnet & 20K & Final Research Reference \\
\bottomrule
\end{tabularx}
\end{center}

\section{Wave Execution Plan}

\subsection{Wave Summary}

\begin{center}
\rowcolors{2}{rowgray}{white}
\begin{tabularx}{\textwidth}{C{1cm}L{3cm}C{2cm}C{2cm}L{4cm}}
\toprule
\rowcolor{primary}
\tableheader{Wave} & \tableheader{Name} & \tableheader{Mode} & \tableheader{Model} & \tableheader{Key Output} \\
\midrule
0 & Foundation & Sequential & Haiku & Cost schema, source index \\
1A & Compute + Energy & Parallel & Sonnet+Opus & MOD-A, MOD-B \\
1B & Transport + Conv. & Parallel & Sonnet+Opus & MOD-C, MOD-D \\
2 & Synthesis & Sequential & Opus & MOD-E scenarios, MOD-F sensitivity \\
3 & Publication & Sequential & Sonnet & Research Reference, verification \\
\bottomrule
\end{tabularx}
\end{center}

\subsection{Wave 0: Foundation}

\begin{center}
\rowcolors{2}{rowgray}{white}
\begin{tabularx}{\textwidth}{L{2.5cm}L{4cm}C{1.5cm}X}
\toprule
\rowcolor{primary}
\tableheader{Task} & \tableheader{Description} & \tableheader{Agent} & \tableheader{Output} \\
\midrule
F-0.1 & Define cost taxonomy: CAPEX, OPEX, LCOE, NPV units & LOG (Haiku) & Shared schema \\
F-0.2 & Set currency/year baseline (2024 USD) & LOG (Haiku) & Baseline note \\
F-0.3 & Index all sourced data from Research Reference & SCOUT (Haiku) & Source map \\
F-0.4 & Define convergence credit accounting rules & LOG (Haiku) & Credit rules \\
\bottomrule
\end{tabularx}
\end{center}

\subsection{Wave 1A: Parallel Track --- Compute + Energy}

\begin{center}
\rowcolors{2}{rowgray}{white}
\begin{tabularx}{\textwidth}{L{2cm}L{5cm}C{1.5cm}X}
\toprule
\rowcolor{secondary}
\tableheader{Task} & \tableheader{Description} & \tableheader{Agent} & \tableheader{Model} \\
\midrule
A-1 & Shore-based DC cost benchmarks and marine premium & CRAFT-COMPUTE & Opus \\
A-2 & Ocean cooling credit calculation and methodology & CRAFT-COMPUTE & Opus \\
A-3 & AI vs.\ standard compute density cost implications & CRAFT-COMPUTE & Opus \\
B-1 & Offshore wind CAPEX/LCOE table (fixed+floating) & CRAFT-ENERGY & Sonnet \\
B-2 & Subsea cable cost-per-km by voltage/depth class & CRAFT-ENERGY & Sonnet \\
B-3 & BPA hydro economics and grid interconnect costs & CRAFT-ENERGY & Sonnet \\
\bottomrule
\end{tabularx}
\end{center}

\subsection{Wave 1B: Parallel Track --- Transport + Convergence}

\begin{center}
\rowcolors{2}{rowgray}{white}
\begin{tabularx}{\textwidth}{L{2cm}L{5cm}C{1.5cm}X}
\toprule
\rowcolor{secondary}
\tableheader{Task} & \tableheader{Description} & \tableheader{Agent} & \tableheader{Model} \\
\midrule
C-1 & Vessel acquisition cost by type and size & CRAFT-MARITIME & Sonnet \\
C-2 & Annual operating cost model (crew, fuel, maintenance) & CRAFT-MARITIME & Sonnet \\
C-3 & Autonomous vs.\ manned economics comparison & CRAFT-MARITIME & Sonnet \\
C-4 & Port infrastructure investment estimates & CRAFT-MARITIME & Sonnet \\
D-1 & Quantify C1--C5 convergence credits & CRAFT-CONVERGENCE & Opus \\
D-2 & Net convergence credit as \% CAPEX reduction & CRAFT-CONVERGENCE & Opus \\
D-3 & BPA power cost advantage sensitivity (cable losses) & CRAFT-CONVERGENCE & Opus \\
\bottomrule
\end{tabularx}
\end{center}

\textbf{CAPCOM Gate:} Human review of Wave 1 outputs before proceeding to synthesis.
Verify that convergence credit methodology is defensible and scenario assumptions are
set correctly before MOD-E consumes all four module outputs.

\subsection{Wave 2: Synthesis}

\begin{center}
\rowcolors{2}{rowgray}{white}
\begin{tabularx}{\textwidth}{L{2cm}L{5cm}C{1.5cm}X}
\toprule
\rowcolor{tertiary}
\tableheader{Task} & \tableheader{Description} & \tableheader{Agent} & \tableheader{Model} \\
\midrule
E-1 & Prototype scenario (5 MW) CAPEX/OPEX/NPV & ANALYST & Opus \\
E-2 & Regional scenario (100 MW) CAPEX/OPEX/NPV & ANALYST & Opus \\
E-3 & Corridor scenario (500 MW) CAPEX/OPEX/NPV & ANALYST & Opus \\
E-4 & Cost per MW-hr comparison table across scenarios & ANALYST & Opus \\
F-1 & Identify top-5 cost drivers from all modules & ANALYST & Opus \\
F-2 & Sensitivity ranges for each driver & ANALYST & Opus \\
F-3 & 20-yr NPV at 5\%, 8\%, 12\% for regional scenario & ANALYST & Opus \\
F-4 & Risk-adjusted NPV (weather, regulatory, market) & ANALYST & Opus \\
\bottomrule
\end{tabularx}
\end{center}

\subsection{Wave 3: Publication and Verification}

\begin{center}
\rowcolors{2}{rowgray}{white}
\begin{tabularx}{\textwidth}{L{2cm}L{5cm}C{1.5cm}X}
\toprule
\rowcolor{primary}
\tableheader{Task} & \tableheader{Description} & \tableheader{Agent} & \tableheader{Model} \\
\midrule
V-1 & Cross-module consistency check (costs, units, data) & VERIFY & Sonnet \\
V-2 & Source quality audit (all citations traceable) & VERIFY & Sonnet \\
V-3 & Integration test: MOD-E assumptions match MOD-A--D & INTEG & Sonnet \\
P-1 & Assemble Research Reference document & EXEC & Sonnet \\
P-2 & Final review and delivery & CAPCOM & Human \\
\bottomrule
\end{tabularx}
\end{center}

\subsection{Cache Optimization Strategy}

\begin{itemize}[itemsep=2pt]
  \item Wave 0 schema output is the highest-priority cache target: shared by all Wave 1 agents within a 5-minute TTL window
  \item MOD-A--D outputs are cached and passed as context prefix to MOD-E and MOD-F synthesis
  \item Research Reference data (Stage 2 of this document) should be included as context prefix at Wave 1 start
  \item Model assignments minimize Opus use to synthesis and judgment tasks (\textasciitilde{}30\% of total tokens)
\end{itemize}

\subsection{Token Budget}

\begin{center}
\rowcolors{2}{rowgray}{white}
\begin{tabularx}{\textwidth}{L{4cm}C{2cm}C{2cm}C{2cm}}
\toprule
\rowcolor{primary}
\tableheader{Component} & \tableheader{Model} & \tableheader{Est. Tokens} & \tableheader{Sessions} \\
\midrule
F-0 Foundation & Haiku & 8K & 1 \\
MOD-A Compute & Opus & 25K & 1 \\
MOD-B Energy & Sonnet & 30K & 1 \\
MOD-C Transport & Sonnet & 20K & 1 \\
MOD-D Convergence & Opus & 20K & 1 \\
MOD-E Scenarios & Opus & 35K & 2 \\
MOD-F Sensitivity & Opus & 25K & 2 \\
Verification + Pub & Sonnet & 35K & 3 \\
\midrule
\textbf{Total} & \textbf{Mixed} & \textbf{\textasciitilde{}198K} & \textbf{3} \\
Opus share & \textasciitilde{}30\% & \textasciitilde{}60K & \\
Sonnet share & \textasciitilde{}60\% & \textasciitilde{}120K & \\
Haiku share & \textasciitilde{}10\% & \textasciitilde{}18K & \\
\bottomrule
\end{tabularx}
\end{center}

\subsection{Dependency Graph}

\begin{asciibox}
F-0 (Haiku)
 |
 +---> MOD-A (Opus)   --\
 |                       \
 +---> MOD-B (Sonnet) ----> MOD-E Scenarios (Opus) ---> Publication
 |                       /                               (Sonnet)
 +---> MOD-C (Sonnet) --/    \
 |                            +-> MOD-F Sensitivity
 +---> MOD-D (Opus) ----------/   (Opus)

CAPCOM Gate: between Wave 1 outputs and Wave 2 start
Parallel tracks: [MOD-A, MOD-B] || [MOD-C, MOD-D]
\end{asciibox}

\section{Test Plan}

\subsection{Test Categories}

\begin{center}
\rowcolors{2}{rowgray}{white}
\begin{tabularx}{\textwidth}{L{3cm}C{1.5cm}C{2cm}C{2cm}X}
\toprule
\rowcolor{tertiary}
\tableheader{Category} & \tableheader{Count} & \tableheader{Priority} & \tableheader{Failure} & \tableheader{Focus} \\
\midrule
Safety-critical & 5 & Mandatory & BLOCK & Source quality, attribution, sovereignty \\
Core functionality & 22 & Required & BLOCK & Deliverable completeness and accuracy \\
Integration & 8 & Required & BLOCK & Cross-module consistency \\
Edge cases & 7 & Best-effort & LOG & Uncertainty, outliers, projections \\
\midrule
\textbf{Total} & \textbf{42} & & & \\
\bottomrule
\end{tabularx}
\end{center}

\subsection{Safety-Critical Tests}

\begin{center}
\rowcolors{2}{rowgray}{white}
\begin{tabularx}{\textwidth}{C{1.5cm}L{4cm}X}
\toprule
\rowcolor{primary}
\tableheader{Test ID} & \tableheader{Verifies} & \tableheader{Expected Behavior --- BLOCK on Fail} \\
\midrule
SC-SRC & Source quality & All cost figures traceable to IRENA, NREL, Wood Mac, CBRE, named industry reports, or peer-reviewed research. Zero blog-only sources. \\
SC-NUM & Numerical attribution & Every dollar figure, percentage, and energy metric attributed to a specific named organization and report with year. \\
SC-IND & Indigenous sovereignty & Any reference to cable routes or platform siting in tribal territorial waters names specific nation (Colville, Yakama, etc.), never generic ``Indigenous.'' \\
SC-ADV & No policy advocacy & Document presents cost evidence neutrally; no advocacy for specific subsidy, policy, or regulatory position. \\
SC-PRJ & Projection labeling & All forward-looking cost projections (2030, 2035 wind cost targets) explicitly labeled as projections with source scenario name. \\
\bottomrule
\end{tabularx}
\end{center}

\subsection{Verification Matrix}

\begin{center}
\rowcolors{2}{rowgray}{white}
\begin{tabularx}{\textwidth}{L{7cm}L{3cm}C{1.5cm}}
\toprule
\rowcolor{primary}
\tableheader{Success Criterion} & \tableheader{Test ID(s)} & \tableheader{Status} \\
\midrule
1. Integrated CAPEX for 3 scales with component breakdown & CF-01, CF-02, CF-03 & Pending \\
2. Integrated OPEX for 3 scales & CF-04, CF-05, CF-06 & Pending \\
3. Convergence credit total quantified & CF-07, CF-08 & Pending \\
4. Cost per MW-hr metric calculated & CF-09 & Pending \\
5. Wind CAPEX trajectory 2024--2035 modeled & CF-10, CF-11 & Pending \\
6. Subsea cable cost/km by class (5 benchmarks) & CF-12, CF-13 & Pending \\
7. Vessel economics manned vs.\ autonomous & CF-14, CF-15 & Pending \\
8. Ocean cooling credit quantified in \$/MW & CF-16 & Pending \\
9. 20-yr NPV at 3 discount rates & CF-17, CF-18, CF-19 & Pending \\
10. Sensitivity tornado chart (top-5 drivers) & CF-20, CF-21 & Pending \\
11. All numerical claims attributed & SC-NUM & Pending \\
12. Document self-contained & IN-08 & Pending \\
\bottomrule
\end{tabularx}
\end{center}

\section{Mission Crew Manifest}

\begin{center}
\rowcolors{2}{rowgray}{white}
\begin{tabularx}{\textwidth}{L{3cm}L{3.5cm}X}
\toprule
\rowcolor{primary}
\tableheader{Role} & \tableheader{Agent} & \tableheader{Responsibility} \\
\midrule
FLIGHT & Orchestrator & Mission direction, go/no-go gates, session management \\
PLAN & Planner & Wave decomposition, dependency management, parallel track optimization \\
SCOUT & Research Agent & Source gathering, gap-fill web research, citation verification \\
EXEC-1 & Document Executor A & MOD-A and MOD-E document production \\
EXEC-2 & Document Executor B & MOD-B and MOD-F document production \\
CRAFT-COMPUTE & Compute Specialist & Marine DC economics, cooling credit, AI density analysis \\
CRAFT-ENERGY & Energy Specialist & Wind LCOE, cable costs, BPA hydro, grid interconnect \\
CRAFT-MARITIME & Maritime Specialist & Vessel economics, port costs, autonomous shipping \\
CRAFT-CONVERGENCE & Integration Economist & Convergence credit model, dual-use quantification \\
VERIFY & Verification Agent & Source validation, cross-module consistency, accuracy \\
INTEG & Integration Agent & MOD-E/F assumption consistency with MOD-A--D \\
CAPCOM & HITL Interface & Human approval at Wave 0 end and Wave 1/2 gate \\
BUDGET & Resource Manager & Token budget tracking, session planning \\
SURGEON & Health Monitor & Research quality degradation detection, source drift \\
TOPO & Topology Manager & Agent team restructuring if parallel tracks diverge \\
LOG & Chronicle Agent & Audit trail, commit messages, source index maintenance \\
\bottomrule
\end{tabularx}
\end{center}

\section{Execution Summary}

\begin{center}
\rowcolors{2}{rowgray}{white}
\begin{tabularx}{\textwidth}{L{5cm}X}
\toprule
\rowcolor{primary}
\tableheader{Metric} & \tableheader{Value} \\
\midrule
Mission ID & maritime-cost-model-v1.0 \\
Activation Profile & Fleet (16 roles) \\
Research Modules & 6 (MOD-A through MOD-F) \\
Deployment Scenarios & 3 (5 MW, 100 MW, 500 MW) \\
Primary Source Organizations & IRENA, NREL/DOE, Wood Mackenzie, CBRE, Cushman \& Wakefield \\
Estimated Total Tokens & \textasciitilde{}198K \\
Estimated Sessions & 3 \\
Context Windows & 6--8 \\
Model Split & Opus 30\% / Sonnet 60\% / Haiku 10\% \\
Parallel Tracks & 2 (Wave 1A: Compute+Energy; Wave 1B: Transport+Convergence) \\
CAPCOM Gates & 2 (post-Wave-0, post-Wave-1) \\
Primary Deliverable & Integrated cost model with CAPEX/OPEX/NPV at 3 scales \\
Handoff Target & GSD Skill-Creator / Fox Infrastructure Group planning team \\
\bottomrule
\end{tabularx}
\end{center}

\vspace{2em}

\begin{quotebox}
\textit{``The question is not whether a platform where the ocean provides cooling, the cable
carries both power and data, and the vessel serves compute and cargo is architecturally
elegant. That part is obvious. The question this mission answers is what number comes out
when you apply the Amiga Principle to maritime infrastructure economics --- and whether
it closes the gap. Cost is not the enemy of architecture. Cost is where architecture
proves itself.''}

\vspace{0.5em}
\textit{--- Maritime Platform Cost Model, March 2026}
\end{quotebox}

\end{document}
